% =========================================================
% Limits - Definitions
% =========================================================
\subsection*{Definitions}

\begin{definition*}[Intuitive Limit Idea]
The limit \(\lim_{x\to a} f(x)=L\) says that the values of \(f(x)\) can be
made close to \(L\) by taking \(x\) close to \(a\), without requiring
\(f(a)\) itself to be defined or equal to \(L\).
\end{definition*}

\begin{definition*}[Epsilon-Delta Reference Definition]
The statement \(\lim_{x\to a} f(x)=L\) means that for every
\(\varepsilon>0\) there is a \(\delta>0\) such that
\[
  0<|x-a|<\delta
  \quad\Longrightarrow\quad
  |f(x)-L|<\varepsilon.
\]
In Volume VI this definition is mainly a reference point for interpreting
limit notation; detailed proofs belong in Real Analysis.
\end{definition*}

\begin{definition*}[One-Sided Limits]
The right-hand limit \(\lim_{x\to a^+} f(x)\) tracks values with \(x>a\).
The left-hand limit \(\lim_{x\to a^-} f(x)\) tracks values with \(x<a\).
The two-sided limit exists exactly when both one-sided limits exist and agree.
\end{definition*}

\begin{definition*}[Limits at Infinity]
The notation \(\lim_{x\to\infty} f(x)=L\) describes the value approached by
\(f(x)\) as \(x\) becomes arbitrarily large. The notation
\(\lim_{x\to a} f(x)=\infty\) describes unbounded growth near \(a\).
\end{definition*}

\begin{definition*}[Indeterminate Forms]
Forms such as \(0/0\), \(\infty/\infty\), \(0\cdot\infty\),
\(\infty-\infty\), \(1^\infty\), \(0^0\), and \(\infty^0\) do not determine
a limit by inspection. They signal that algebraic or analytic rewriting is
needed.
\end{definition*}

\begin{definition*}[Asymptote]
A horizontal, vertical, or oblique asymptote records limiting behavior of a
graph near infinity or near a point where the function grows without bound.
\end{definition*}
